\documentclass[12pt]{article}
\usepackage[margin=1in]{geometry}
\usepackage{setspace}
\usepackage{parskip}
\usepackage{titlesec}
\usepackage{enumitem}
\usepackage{booktabs}
\usepackage{array}
\usepackage{hyperref}
\usepackage{fancyhdr}
\usepackage{graphicx}
\usepackage{xcolor}
\usepackage{listings}
\usepackage{caption}
\usepackage{float}
\usepackage{amssymb}
\usepackage{amsmath}
\usepackage{siunitx}
\usepackage{csquotes}
\usepackage{etoolbox}

%--- Header / Footer ---------------------------------------------------------
\pagestyle{fancy}
\fancyhf{}
\rhead{\textit{The Salt Coin of the Drowned Prince}}
\lhead{\small\textit{A Zakov Expansion for Fate's Edge}}
\rfoot{\thepage}
\renewcommand{\headrulewidth}{0.4pt}
\renewcommand{\footrulewidth}{0.4pt}

%--- Section formatting -------------------------------------------------------
\titleformat{\section}
  {\large\bfseries\sffamily}
  {}{0pt}{}
\titleformat{\subsection}
  {\normalsize\bfseries\sffamily}
  {}{0pt}{}
\titleformat{\subsubsection}
  {\small\bfseries\sffamily}
  {}{0pt}{}
\renewcommand{\sectionmark}[1]{%
  \markright{\thesection\ #1}}
\renewcommand{\subsectionmark}[1]{%
  \markright{\thesubsection\ #1}}

%--- Custom environments ------------------------------------------------------
\newenvironment{quoteBox}
  {\begin{quote}\itshape}
  {\end{quote}}

\newenvironment{mechanic}
  {\begin{list}{}{\setlength{\leftmargin}{1.5em}\setlength{\rightmargin}{0pt}}
   \item[]\textbf{Mechanic:}}
  {\end{list}}

\newenvironment{statblock}
  {\begin{tabular}{@{}l l@{}}}
  {\end{tabular}}

%--- Document -----------------------------------------------------------------
\begin{document}

\begin{center}
    {\LARGE\bfseries The Salt Coin of the Drowned Prince}\\[4pt]
    {\large\itshape A Zakov Expansion for Fate's Edge}\\[12pt]
    \textit{``The tide does not care if you are laughing. It will still drag you down.''}\\[24pt]
    {\large As told by a Sidhi deckhand with a hook for a hand and a laugh like breaking glass.}
\end{center}
\vspace{1em}

\section*{The Story: The Merry Crew of the \textit{Laughing Eel}}
\addcontentsline{toc}{section}{The Story: The Merry Crew of the \textit{Laughing Eel}}

You want a story about Zakov? Fine. But don't say I didn't warn you.

Three seasons back I shipped with Captain Mira “the Lucky” aboard the \textit{Laughing Eel}. She was a Sidhi woman with a smile that could sell sand to a djinn and a debt to every dockmaster in Port Shed. Her crew? A dozen misfits, deserters, and dreamers who could not hold a real job because they were too busy pretending the world was funny.

We laughed. Gods, we laughed. We laughed when the anchor chain snapped. We laughed when the cook's stew grew legs. We laughed when the Salt Prince's tax collector came aboard and we fed him expired rum until he signed our manifest as “inspection passed.” 

It was merry. It was stupid. It was the best time I ever had.

Then we found the coin.

\subsection*{The Find}
It was a calm night in the Drowned Quarter—part of Zakov where the tide covers the streets twice a day and the buildings lean like drunk uncles. We were offloading a cargo of stolen Thepyrgosi glassware (don’t ask) when the eel‑girl, Fen, dove off the pier to cool off.

She came up screaming.

Not in pain. In \textit{wonder}.

“The bottom of the slip is all silver,” she said. “Coins. Hundreds. Maybe thousands.”

We dove. It was true. The seafloor glittered like a madman’s dream. We filled our pockets, our boots, our mouths. We laughed the whole time. Water in our lungs, salt in our eyes, silver in our fists.

Captain Mira held up one coin—big as a palm, stamped with a face that was half skull, half king. “This,” she said, “is going to buy us a new ship.”

The coin was warm. None of us noticed.

\subsection*{The Turn}
That night Fen dreamed of drowning. She woke with brine in her mouth and a single word on her lips: \textit{Oblivion}.

The next day Old Kes—cook, veteran of three shipwrecks and a shark attack—walked off the dock at high tide. He did not jump. He did not trip. He simply stepped, as if the water were a staircase he had always meant to take.

We pulled him out. He was smiling. “The prince wants his silver back,” he said. Then he laughed. It was not his laugh.

That night the \textit{Laughing Eel} sailed itself out of the harbor. We woke to an empty berth and a single silver coin on the dock, warm to the touch.

Captain Mira picked it up. “We’re going after it,” she said.

No one argued. Not because we were brave. Because the coin in her palm was whispering our names.

\subsection*{The Voyage}
The sea turned strange. Waters that should have been green ran black. Fish swam in circles. Gulls stood on the rigging and stared at us with eyes that were too human.

We laughed anyway—nervous laughter, the kind you use to keep the dark from hearing you think.

Fen stopped eating. She said the food tasted like memory. She began humming a tune none of us recognised. It sounded like waves on a shore where no land existed.

Captain Mira called a meeting. “We’re cursed,” she said. “Anyone want to leave?”

No one did. Not because we were loyal. Because none of us could remember the way back to port. The map had gone blank. The stars had moved.

\subsection*{The Prince}
We found the \textit{Laughing Eel} three days later, drifting off a reef that was not on any chart. The crew was gone. The deck was scrubbed clean. In the captain’s quarters, a single silver coin sat on the table.

Beneath it, a note written in water: \textit{“You took what was mine. Now I will take what is yours.”}

That night the Drowned Prince came aboard. He wore a crown of barnacles and a smile that had too many teeth. His voice was the sound of a ship breaking on rocks—slow, patient, and final.

“You took my salt,” he said. “My coin. My silence. I have waited three hundred years for someone to remember me.”

Captain Mira stepped forward. She was shaking. She was also still Captain Mira the Lucky.

“We’ll give it back,” she said. “All of it. With interest.”

The prince laughed. It was the most terrible sound I have ever heard.

“Interest,” he said. “Yes. I like that.”

\subsection*{The Bargain}
He let us live. On one condition: every year, on the same night, we must return to this reef and throw one silver coin into the water—not just any coin, but a coin that has been paid for a life. A coin that was used to buy a person's freedom. A coin that has seen blood.

If we forget, he will come again. And next time, he will not ask for interest.

We have kept our bargain. Every year we find a coin. Every year we return. Every year the sea is calm.

But Fen—Fen, who dreamed of drowning, who hummed the nameless tune—Fen went back alone last spring. She took a boat out at midnight. She did not return.

The coin on her pillow was warm.

She had paid the interest for all of us.

Now we laugh quieter. Now we check the water before we sail. Now we remember that Zakov is not a city of merry pirates. It is a city of debts that live under the waves, waiting for you to forget.

\section*{Mechanics: The Salt Prince’s Drowned Court}
\label{sec:mechanics}

\subsection{New Faction: The Drowned Prince (Terrestrial Patron)}
The Drowned Prince is not a god. He is a \textit{creditor}. He was once a Zakov pirate king who made a bargain with Khemesh (the Abyssal Maw) to never die. Khemesh gave him immortality—but not in the way he expected. His body sank; his spirit rose. Now he rules the reefs, the wrecks, and the coins that have tasted death.

\subsubsection{Gaining His Favor}
Find a salt coin (see below) and return it to his reef. He will offer a bargain: power in exchange for an annual payment.

\subsubsection{Obligation Track}
The Drowned Prince uses a standard Obligation track (Spirit + Presence). However, Obligation to him cannot be cleared by Acts of Service—only by paying the annual toll (1 silver coin per segment owed). If you cannot pay, the sea takes something from you: a memory, a skill point, or a year of your life.

\subsubsection{New Rites (for Runekeepers and Witches)}
\begin{description}[leftmargin=!,labelwidth=\widthof{\bfseries High}]
  \item[Salt Coin’s Toll (Low, 4 XP)] Touch a coin that has been used to buy a life (ransom, slave’s purchase, bribe for execution). For one scene you may reroll one failed roll involving deception, escape, or bargains. Mark +1 Obligation to the Drowned Prince.
  \item[Drowned Prince’s Favor (Standard, 8 XP)] Once per journey you may call on the prince to calm a storm, part a reef, or guide your ship through treacherous waters. Mark +2 Obligation. The GM gains 1 Story Beat (the sea remembers your debt).
  \item[The Interest Comes Due (High, 12 XP)] Once per campaign you may invoke this rite when you are about to die at sea. The Drowned Prince saves you—but you owe him your life. Mark +4 Obligation. You must pay 1 silver coin per segment within a year, or the prince will claim a loved one instead.
\end{description}

\subsection{New Mechanic: Salt Coins}
Salt coins are silver pieces that have been used to pay for a life. They are physically indistinguishable from ordinary coins, but they are always warm. A salt coin can be:
\begin{itemize}
  \item \textbf{Spent as a Boon} (one use) to reroll any failed roll involving the sea, drowning, or escape.
  \item \textbf{Traded to the Drowned Prince} to clear 1 segment of Obligation.
  \item \textbf{Used as material for the Rite of Salt Coin’s Toll.}
\end{itemize}

\textbf{Finding Salt Coins:} On any Miss involving a death, a ransom, or a broken oath, the GM may place a salt coin in the scene. A player who searches (Wits + Investigation vs DV 3) may find it.

\textbf{Curse of the Salt Coin:} A character who carries more than 3 salt coins at once begins to dream of drowning. Each night test Spirit + Resolve (DV 4) or mark 1 Fatigue. Three failed tests in a row and the Drowned Prince sends a \textit{Salt Walker} (see below) to collect.

\subsection{New Creature: Salt Walker (Tier II, Cap 3)}
A drowned sailor raised by the Drowned Prince. It moves slowly, relentlessly, and smells of brine and old regret. It cannot be reasoned with—only bargained with.

\begin{center}
\begin{tabular}{@{}ll@{}}
\toprule
\textbf{Stat} & \textbf{Value} \\
\midrule
Body & 3 \\
Wits & 1 \\
Spirit & 2 \\
Presence & 1 \\
Harm & 4 \\
Fatigue (Body) & 3 \\
\bottomrule
\end{tabular}
\end{center}

\textbf{Traits:}
\begin{itemize}
  \item \textbf{Brine Aura:} Any living creature within Near range suffers –1 die to Endurance rolls (salt burns the lungs).
  \item \textbf{Inexorable:} Salt Walkers cannot be stopped by mundane means. They will walk through walls, across water, and over broken ground. They take 1 Harm per round from fire, but otherwise regenerate 1 Harm per round.
  \item \textbf{Bargain:} Instead of fighting, a character may offer a salt coin to the Walker. The Walker takes the coin and leaves. If no coin is offered, the Walker attacks.
\end{itemize}

\textbf{Resolution:} Salt Walkers are not destroyed; they are \textit{paid}.

\subsection{New Talent: Salt‑Scarred (Major, 8 XP)}
You have survived an encounter with the Drowned Prince, and the sea has marked you.

\begin{itemize}
  \item You can hold your breath for twice as long as normal (test Spirit + Athletics vs DV 3 to extend to 1 minute per success).
  \item Once per session, when you fail a roll involving the sea (swimming, sailing, navigation), you may reroll one die. The GM gains 1 Story Beat (the sea remembers).
  \item You always know which way is north when you can see the horizon. This is instinct, not magic.
  \item \textbf{Drawback:} You cannot sleep within a mile of a salt coin without marking 1 Fatigue. The prince whispers to you in dreams.
\end{itemize}

\subsection{New Mechanic: The Tide Clock (Nautical Pressure)}
For any scene at sea in Zakov waters (or near the Drowned Prince’s reef), the GM may start a \textbf{Tide Clock [6]}. It represents the rising of the prince’s attention.

\begin{itemize}
  \item Tick the timer when:
    \begin{itemize}
      \item A character fails a roll involving salt water.
      \item A salt coin is used or found.
      \item The party lingers in one place for more than an hour.
    \end{itemize}
  \item When the timer reaches 3, the water becomes choppy. All sailing rolls are at –1 die.
  \item When the timer reaches 6, the Drowned Prince sends a Salt Walker (or his attention manifests as a rogue wave, a crack in the hull, or a sudden fog). The GM chooses.
\end{itemize}

\textbf{Clearing the Tide Clock:} Paying a salt coin to the Drowned Prince (via the Rite or a bargain) clears 2 segments. Leaving Zakov waters resets the timer.

\subsection{Adventure Seed: The Interest Comes Due}
The crew of the \textit{Laughing Eel} has missed their annual payment. Fen has been gone for a year. Captain Mira hires the party to find Fen—or her replacement. The Drowned Prince will accept any soul, as long as it is willing. The party must sail to the reef, negotiate with the prince, and decide: pay with a life, pay with a memory, or fight the tide itself.

\begin{description}
  \item[Prince’s Patience [6]] Ticks each day the party delays. When full, the prince sends a Salt Walker.
  \item[Fen’s Breath [4]] Fen is still alive, trapped in the prince’s reef. Each day, the timer ticks. When full, she becomes a Salt Walker herself.
\end{description}

\textbf{Resolution Options:}
\begin{itemize}
  \item \textbf{Pay with a memory:} A character offers a cherished memory (erase a Bond or a positive trait). The prince accepts. Fen is released. The character gains the \textit{Salt‑Scarred} talent for free—but the prince will remember their face.
  \item \textbf{Pay with a life:} The party finds a willing volunteer (a condemned prisoner, a dying elder). The prince takes them. Fen is released. The party gains a Debt to the prince’s cult.
  \item \textbf{Fight the tide:} The party attacks the prince’s physical form (Tier IV, Cap 5). On success, Fen is freed and the prince retreats for a year. On failure, a party member is taken instead.
\end{itemize}

\subsection{New Patron: The Drowned Prince (for Invokers and Witches)}
\begin{itemize}
  \item \textbf{Symbol:} A silver coin worn around the neck, stamped with a skull and crown.
  \item \textbf{Favored Deeds:} Paying a debt, ransoming a prisoner, burying a drowned sailor with respect, throwing a coin to a beggar (the prince sees all silver as salt).
  \item \textbf{Taboos:} Spending a salt coin on frivolity, ignoring a drowning person, leaving a body unburied at sea.
\end{itemize}

\subsubsection{Low Rite – The Salt Oath}
Swear an oath while holding a salt coin. If you break it, the coin turns black and the prince sends a Salt Walker after you. (Mark +1 Obligation.)

\subsubsection{Standard Rite – The Drowned Council}
Once per session you may ask the prince one question about a person who died at sea. He answers truthfully, but cryptically, and the GM gains 2 Story Beats.

\subsubsection{High Rite – The Tide of Oblivion}
You may call upon the prince to flood a zone (Near range). Enemies test Body + Athletics (DV 5) or be swept away. Costs 3 Obligation and 1 salt coin.

\subsection{The Drowned Prince’s Corruption Table (for Cantors and Witches)}
\begin{center}
\begin{tabular}{>{\bfseries}l l l}
\toprule
Tier & Benefit & Cost / Quirk \\
\midrule
1 & Brine‑Sense: +1 die to detect salt water or hidden coins. & Salty Tongue: You taste salt constantly. Social rolls involving food or drink are at –1 die. \\
2 & Tide‑Step: Once per scene, move across water as if it were solid ground. & Drowned’s Memory: You have nightmares of drowning. Mark 1 Fatigue if you don’t sleep near salt water. \\
3 & Prince’s Favor: Once per session, force a reroll on any roll involving the sea or boats. & Salt‑Stained: Your skin becomes rough and grey. –1 die to Persuasion with land‑dwellers. \\
4 & Oblivion’s Touch: Your melee attacks deal +1 Harm against wet or damp targets. & Waking Dread: When you hear waves, you must test Resolve (DV 3) or lose your next action. \\
5 & Prince’s Voice: Once per session, command a single salt‑water creature (fish, eel, crab) to deliver a message. & Debt of the Deep: You owe the prince a favour. He will collect when you least expect it. \\
6+ & Drowned Ascendancy: Once per campaign, you may sink a ship with a touch. The prince claims every soul. & The Interest Never Ends: Mark +3 Obligation permanently. You can never fully leave the prince’s service. \\
\bottomrule
\end{tabular}
\end{center}

\section*{Final Note from the Narrator}
\begin{quoteBox}
Zakov is not a city of friendly pirates. It is a city of debts that breathe salt and speak in whispers. The \textit{Laughing Eel} still sails, but the crew laughs quieter now. They check the water before they drink. They never go to the reef alone.

And if you ever find a silver coin that is warm to the touch? Spend it fast. Or better yet, throw it back.

The prince is patient. He can wait for the tide.

– A deckhand who knows better, Zakov docks, year of the broken anchor
\end{quoteBox}

\end{document}
